\section{Pappus's hexagon theorem}

\subsection{Statement}

Let $A,B,C$ be three points on one projective line and $D,E,F$ three points on
another.  Suppose the two base lines are distinct and all indicated joins and
meets are defined.  Put
\[
 X=AE\cap BD,
 \qquad Y=AF\cap CD,
 \qquad Z=BF\cap CE.
\]
Then $X,Y,Z$ are collinear.

\subsection{Constructive projective normalization}

Let $O$ be the intersection of the two base lines.  Choose one of the named
points on each line away from $O$; the three resulting noncollinear vectors
form a projective basis.  A fixed finite case split selects such points under
the stated nondegeneracy hypotheses.  In this basis the lines are $z=0$ and
$y=0$, and we may write
\[
 A=(a,1,0),\quad B=(b,1,0),\quad C=(c,1,0),
\]
\[
 D=(d,0,1),\quad E=(e,0,1),\quad F=(f,0,1).
\]
The change of basis is an explicit adjugate divided by a nonzero determinant,
so this normalization is available constructively and is primitive recursive
on rational codes.

\subsection{The three intersections}

Cross products give
\begin{align*}
 X&=(A\join E)\meet(B\join D)
   =\bigl(ad-be,\ d-e,\ a-b\bigr),\\
 Y&=(A\join F)\meet(C\join D)
   =\bigl(ad-cf,\ d-f,\ a-c\bigr),\\
 Z&=(B\join F)\meet(C\join E)
   =\bigl(be-cf,\ e-f,\ b-c\bigr).
\end{align*}
The decisive certificate is not even a determinant expansion:
\begin{equation}\label{eq:pappus-linear-relation}
 X-Y+Z=0.
\end{equation}
Indeed, each of the three coordinates cancels term by term.  Equation
\eqref{eq:pappus-linear-relation} makes $X,Y,Z$ linearly dependent; therefore
\[
 \det(X,Y,Z)=0,
\]
which is exactly their collinearity.

If, for example, $X$ and $Y$ are distinct, the explicit line witness is
$\ell=X\cross Y$.  The first two incidences are cross-product identities and
\[
 \ell\cdot Z=\det(X,Y,Z)=0.
\]
Thus the proof also supplies the line containing the three points.

\subsection{Degeneracies and the polynomial closure}

The vector identity remains true without dividing by any coordinate.  If one
of the displayed intersections is the zero cross product, the usual geometric
point is undefined; that case is excluded by the theorem's incidence
hypotheses.  Coincidences among $X,Y,Z$ cause no difficulty: the determinant
conclusion still holds, and a containing line can be taken from any distinct
pair supplied by the configuration.  The generic proof therefore extends to
all legitimate specializations by the same polynomial identity, not by a
classical continuity argument.
